\section*{Para la operación}

\begin{itemize}
\item Depurar la taxonomía de clases antes de ampliar el alcance de la herramienta. El análisis de errores del modelo identificó pares de clases que designan el mismo concepto y difieren únicamente en puntuación, acentuación o separador. Mientras esas duplicidades persistan, una fracción del error de clasificación es irreducible por cualquier modelo, y los reportes de consumo por categoría seguirán fragmentando información que corresponde a un mismo concepto. La depuración no requiere desarrollo adicional: el propio análisis de confusiones señala los candidatos a fusión, y su corrección mejora simultáneamente la calidad del catálogo y el desempeño medido de la herramienta.

\item Hacer exigible la convención de separadores en el momento de la captura. El análisis de calidad identificó 3,649 descripciones que no siguen la convención de dos puntos o punto y coma, con tipos de material en los que esa proporción supera el treinta por ciento. Una descripción sin esa estructura es indistinguible de otra similar mediante una búsqueda exacta, de modo que cada registro capturado sin separador incorpora al catálogo un riesgo de duplicidad que ninguna consulta posterior detecta con certeza. Validar la estructura en el momento de la captura evita que el pasivo siga creciendo mientras se depura el acumulado, y la plataforma ya realiza esa normalización antes de proponer la descripción al gestor.

\item Resolver la categorización de los materiales sin clase asignada. Los 5,368 registros que carecen de categoría en la fuente quedaron fuera del conjunto de modelado y, por lo tanto, fuera del alcance del clasificador. Constituyen un volumen suficiente para justificar una campaña de categorización asistida por la propia plataforma, que puede procesarlos por lotes aprovechando la sugerencia de categoría y la revisión del gestor.

\item Consolidar la operación en esta línea de negocio antes de escalar a otras unidades. El equipo de gestores del área es reducido, de modo que el volumen atendido durante la validación no refleja una adopción parcial sino el tamaño real de la operación. Lo que corresponde entonces no es ampliar el número de usuarios dentro de la misma unidad, sino sostener el uso durante un período más extenso que permita acumular decisiones suficientes para medir la calidad del catálogo con la misma solidez con que ya se midió el tiempo de gestión. Ese resultado es el criterio con el que debería decidirse la extensión de la herramienta a otras líneas de negocio de la corporación: escalar antes de consolidar trasladaría a un ámbito mayor conclusiones que todavía descansan sobre cinco semanas de operación en una sola unidad.

\item Contrastar el umbral de confianza con el comportamiento real de la operación. El umbral vigente se fijó a partir de la partición de prueba, donde el modelo acierta el 98.9\,\% de los casos que resuelve por sí solo. Conviene verificar que esa promesa se cumple en la práctica, y la plataforma ya registra lo necesario para hacerlo: cada vez que un gestor modifica la categoría sugerida, la corrección queda almacenada. El contraste consiste en comparar la proporción de correcciones sobre las solicitudes que el sistema resolvió automáticamente contra el margen de error que el umbral admite. Si las correcciones resultan más frecuentes de lo previsto, el umbral es demasiado permisivo y conviene elevarlo, a costa de derivar más solicitudes a revisión del gestor; si resultan menos frecuentes, puede reducirse para automatizar un volumen mayor sin comprometer la calidad. Esta revisión debe repetirse después de cada reentrenamiento, porque un modelo nuevo altera la distribución de confianzas sobre la que el umbral fue ajustado.

\item Explotar la instrumentación de calidad que el sistema ya captura. Las vistas del esquema gold registran la tasa de corrección, las decisiones sobre duplicados y el desglose por etapa del proceso, información que durante el período de validación no se explotó. Incorporarla al seguimiento habitual del área convierte el tablero en un instrumento de gestión y no solo en un reporte del proyecto.
\end{itemize}

\section*{Para la continuidad técnica}

\begin{itemize}
\item Adoptar una partición agrupada por descripción en el procedimiento de evaluación. La duplicidad del catálogo permite que registros con descripción idéntica queden repartidos entre entrenamiento y prueba, lo que introduce un sesgo optimista en las métricas reportadas. Sustituir la partición aleatoria por una que asigne todos los registros de una misma descripción a una sola de las particiones corrige el sesgo en la fuente y hace que cada reentrenamiento futuro publique métricas comparables entre versiones.

\item Sustituir la partición única por validación cruzada. Las métricas actuales provienen de una sola partición, de modo que no admiten intervalos de confianza y las diferencias observadas entre configuraciones no pueden declararse estadísticamente distinguibles. Una validación cruzada estratificada permitiría acompañar cada métrica de su dispersión y sostener con evidencia las comparaciones entre algoritmos.

\item Realizar la evaluación por ablación del término principal. La convención de nomenclatura hace que la primera palabra de la descripción coincida con frecuencia con el nombre de la clase. Medir la caída de desempeño al eliminar ese término cuantificaría qué proporción del acierto depende de la correspondencia léxica directa y cuánto de la capacidad discriminativa del modelo.

\item Cuantificar la duplicidad difusa sobre el catálogo completo. El análisis de calidad reportado en este documento mide únicamente duplicidad exacta (texto idéntico); el riesgo de duplicidad que la búsqueda exacta no detecta, descripciones que refieren al mismo material con distinta abreviatura, orden o separador, no se ha cuantificado de forma sistemática. Ejecutar el motor de similitud difusa de la plataforma sobre el maestro completo, al mismo umbral con el que opera en producción, permitiría dimensionar ese riesgo con una metodología trazable y comparable directamente con el comportamiento de la herramienta en operación.

\item Evaluar la incorporación del clasificador al flujo de carga masiva. El trabajo se acotó deliberadamente a la creación individual de materiales. El mismo modelo, aplicado a las cargas masivas que hoy ingresan sin verificación de categoría, extendería el alcance del gobierno de datos sin requerir desarrollo analítico adicional.
\end{itemize}